\documentclass[11pt]{article}
\usepackage[margin=1in]{geometry}
\usepackage{hyperref}
\usepackage{listings}
\usepackage{xcolor}
\usepackage{amsmath}

\lstset{
  basicstyle=\ttfamily\small,
  breaklines=true,
  columns=fullflexible,
  frame=single,
  backgroundcolor=\color{gray!10},
  showstringspaces=false
}

\title{URI OPeNDAP Test Case: Working Log (Prompt \#4)}
\author{Compiled for P. Cornillon}
\date{2026-05-16}

\begin{document}
\maketitle

\section*{Purpose}

Working log for Prompt~\#4 of the URI OPeNDAP test case: estimate
the spatial range over which the per-pixel SST gradients were
computed for the data summarized in
\texttt{gradients\_by\_period/}. Findings are in
\texttt{docs/uri\_test\_case\_4.tex}.

\section{Prompt \#4 (2026-05-16 21:40 UTC)}

\subsection{Instruction}

\emph{``Estimate the spatial range over which the gradients were
calculated for branch \#2 and explain how you arrived at the
answer. This information will be used in the next prompt.''}

\subsection{Interpretation}

The gradient \emph{sums} in \texttt{gradients\_by\_period/} are
1$^\circ$ monthly bins of per-pixel SST gradients computed at the
L2 swath level in the orbit files
(\texttt{SST\_Orbits/}, fields \texttt{eastward\_gradient} and
\texttt{northward\_gradient}). So the \emph{spatial range over
which a single gradient value was calculated} reduces to the
support of the differentiation operator that maps the orbit-level
SST field into one of the orbit-level gradient values. That is
the question I attack below. Independently, of course, the
monthly file aggregates gradients spatially into 1$^\circ$ cells,
but that 1$^\circ$ binning is a \emph{post-hoc} summation, not the
spatial range over which the underlying gradient is computed.

\subsection{Approach}

Three independent estimates, in order of cheapness.

\paragraph{Approach 1: pixel-count ratios from Prompt~\#3.}
The monthly file already tells us something. Per cell:
\texttt{day\_pixel\_count} = 10\,030 valid SST pixels on average;
\texttt{day\_as\_pixel\_count} = 8\,110 valid along-scan gradient
pixels; \texttt{day\_at\_pixel\_count} = 4\,090 valid along-track
gradient pixels. MODIS Aqua's 10-line scan structure means that
along-track derivatives that straddle a scan boundary are masked.
For a centered finite-difference operator with along-track
half-width $H$ pixels, the fraction of valid SST pixels with valid
along-track gradient is approximately
$(10-2H)/10$ minus an additional cloud-edge loss. Subtracting
the cloud-edge loss inferred from along-scan
(\texttt{N\_at}/\texttt{N\_pix} = 0.41,
\texttt{N\_as}/\texttt{N\_pix} = 0.81; so the scan-boundary part
of the along-track deficit is about
$0.81 - 0.41 = 0.40$ which equals $2H/10$), gives $H\approx 2$
pixels, i.e.\ an along-track operator support of $\sim$4 km.

\paragraph{Approach 2: match stored eastward gradient against
centered finite differences of the orbit SST at multiple scales.}
\texttt{scripts/grad\_scale\_findregion.py} walks the
2024-01-01 orbit 115220 for a 1000-line window with the most
valid gradient pixels; the winning window is at $\sim$20\,N,
180\,E (central Pacific, north of Hawaii). I pulled a
$200\times 1354$-pixel block of \texttt{SST\_In},
\texttt{regridded\_sst}, \texttt{eastward\_gradient},
\texttt{northward\_gradient}, \texttt{latitude},
\texttt{longitude}, \texttt{refined\_mask}, and \texttt{qual\_sst}
to \texttt{/tmp/grad\_block2.pkl}. Then
\texttt{scripts/grad\_scale\_magnitude.py} compares stored
$|\nabla T|$ to candidate $|\nabla T|$ computed as centered finite
differences of (a) raw \texttt{SST\_In} and (b)
\texttt{regridded\_sst}, at scales $h\in\{1,2,3,4,5,6,7,8,10\}$
pixels (full operator support $2h$ km using the median across-track
spacing of 1.35 km). Doing this on the magnitude (rather than the
signed component) sidesteps the swath-vs-geographic rotation.
Results:

\begin{itemize}
  \item For \texttt{regridded\_sst} input, the per-pixel RMSE
    between stored and candidate $|\nabla T|$ has a clear minimum
    at $h=3$ (support $\sim$6 km).
  \item For \texttt{SST\_In} input the same RMSE minimum is at
    $h=3$ but with a worse residual; \texttt{regridded\_sst}
    is the operative input field.
  \item Mean $|\nabla T|$ from \texttt{regridded\_sst}-based
    centered diff: matches stored at $h \in [2,3]$ (operator
    support 4--6 km). Mean from \texttt{SST\_In}: matches stored
    at $h\approx 10$ (support $\sim$20 km), reflecting that
    \texttt{SST\_In} is much noisier than the field actually
    used to compute the gradient.
\end{itemize}

The agreement between scales picked by ``minimum RMSE'' and ``mean
magnitude match'' for \texttt{regridded\_sst} as input is the
strongest single piece of evidence: the gradient operator is well
modelled by a centered finite difference with full support
4--6 km, applied to \texttt{regridded\_sst}.

\paragraph{Approach 3: autocorrelation of the stored gradient.}
\texttt{scripts/grad\_scale\_autocorr.py} computes the lagged
correlation of $|\nabla T|$ in the central swath. For a gradient
field computed with operator support $L$ pixels, the
autocorrelation should remain elevated for lags $< L$ (adjacent
pixel pairs share input pixels) and drop sharply beyond $L$.
Observed:
\begin{itemize}
  \item Across-track: AC drops $1 \to 0.64$ at 1.35 km, $\to 0.32$
    at 2.7 km, $\to 0.19$ at 4 km, $\to 0.10$ at 6.75 km, and is
    flat at $\sim$0.05--0.10 beyond 8 km. Effective decorrelation
    scale $\approx$ 2--3 pixels = $\sim$3--4 km.
  \item Along-track: AC drops $1 \to 0.50$ at 1.13 km, $\to 0.18$
    at 2.26 km, $\to 0.13$ at 3.4 km, $\to 0.07$ at 4.5 km, near
    zero by 5--6 km. Decorrelation scale $\approx$ 2 pixels =
    $\sim$2--3 km.
  \item For comparison: the same autocorrelation computed on a
    1-pixel centered difference of \texttt{SST\_In} drops to
    \emph{negative} values at lag 2--3 (Nyquist-like
    anti-correlation of a sharp derivative), returning to zero
    only around 8 km. The stored gradient does not exhibit this
    anti-correlation, ruling out a single 1-pixel diff of the
    raw \texttt{SST\_In}.
\end{itemize}

\subsection{Convergent estimate}

All three approaches agree on a number near
\textbf{$5\pm 2$ km}. The pixel-count approach gives $\sim$4 km;
the magnitude-matching approach gives 4--6 km; the
autocorrelation approach gives 3--4 km across-track and 2--3 km
along-track. The slight anisotropy is consistent with MODIS
Aqua's pixel geometry (1.35 km across-track vs.\ 1.13 km
along-track near nadir) plus the scan-boundary masking which
affects only along-track derivatives.

\subsection{Limitations of the estimate}

\begin{itemize}
  \item The operator may not be a pure centered difference. The
    actual algorithm is in URI's processing code, which I have
    not read. A Sobel, prewitt, or local least-squares operator
    of similar effective support would give similar empirical
    fingerprints.
  \item The \texttt{regridded\_sst} field is itself derived from
    \texttt{SST\_In} by a \texttt{refined\_mask} step that removes
    pixels flagged as ``original high gradient values''. That step
    is a mask, not a spatial average, so its effect on the
    apparent operator support is small. If it also includes a
    box-mean or similar averaging step (which I cannot verify from
    the OPeNDAP-exposed metadata), the total spatial range of
    influence per gradient value would be slightly larger.
  \item The sample region is a tropical Pacific area
    ($\sim$20\,N, 180\,E) without strong fronts. The operator
    support estimate is a statistical-mean property; in regions
    with sharp fronts the spectral content of the gradient field
    differs and an autocorrelation-based estimate would shrink
    somewhat.
\end{itemize}

\subsection{Commands}

\begin{lstlisting}
python3 scripts/grad_scale_findregion.py   # locate best swath block
python3 scripts/grad_scale_block2.py       # pull the block to /tmp
python3 scripts/grad_scale_magnitude.py    # match by magnitude RMSE
python3 scripts/grad_scale_autocorr.py     # autocorrelation analysis
\end{lstlisting}

\subsection{Session timing}

Started 2026-05-16 21:40\,UTC. Closed log 2026-05-16
$\sim$22:15\,UTC.

\end{document}
